\subsection{Morfismos localmente de tipo finito}
\label{subsection:IV.1.3-fr}

\begin{env}[1.3.1]
\label{IV.1.3.1-fr}
Recordemos que si $A$ es un anillo, se dice que una $A$-álgebra (conmutativa) $B$ es
\emph{de tipo finito} si está generada por un número finito de elementos de
$B$, o, lo que es lo mismo, si es isomorfa al cociente de un
álgebra de polinomios $A[T_1,\ldots,T_n]$ por un ideal de esta álgebra. Está
claro que para toda $A$-álgebra (conmutativa) $A'$,
$B\otimes_AA'$ es entonces una $A'$-álgebra de tipo finito. Si $B$ es una
$A$-álgebra de tipo finito y $C$ una $B$-álgebra de tipo finito, $C$ es una
$A$-álgebra de tipo finito, pues si $C$ es un cociente de
$B[T_1,\ldots,T_n]$ y $B$ un cociente de $A[S_1,\ldots,S_m]$,
$B[T_1,\ldots,T_n]=B\otimes_AA[T_1,\ldots,T_n]$ es un cociente de
$A[S_1,\ldots,S_m,T_1,\ldots,T_n]$, por lo que lo mismo ocurre con $C$.
\end{env}

\begin{definition}[1.3.2]
\label{IV.1.3.2-fr}
Sean $f:X\to Y$ un morfismo de preesquemas, $x$ un punto de $X$,
$y=f(x)$. Se dice que $f$ es de tipo finito en $x$ si existe una
vecindad abierta afín $V$ de $y$ y una vecindad abierta afín $U$ de $x$
tales que $f(U)\subset V$ y que el anillo $A(U)$ sea una $A(V)$-álgebra de
tipo finito. Se dice que $f$ es localmente de tipo finito si es de tipo
finito en todo punto de $X$.
\end{definition}

\oldpage[IV-1]{229}

\begin{proposition}[1.3.3]
\label{IV.1.3.3-fr}
Si $Y$ es un preesquema localmente noetheriano y $f:X\to Y$ un morfismo
localmente de tipo finito, entonces $X$ es localmente noetheriano.
\end{proposition}

Para la demostración, véase (\textbf{I}, 6.3.7).

\begin{proposition}[1.3.4]
\label{IV.1.3.4-fr}
\begin{enumerate}
  \item[\textup{(i)}] Toda inmersión local es localmente de tipo finito.
  \item[\textup{(ii)}] Si dos morfismos $f:X\to Y$, $g:Y\to Z$ son localmente de
  tipo finito, lo mismo ocurre con $g\circ f$.
  \item[\textup{(iii)}] Si $f:X\to Y$ es un $S$-morfismo localmente de tipo finito,
  $f_{(S')}:X_{(S')}\to Y_{(S')}$ es localmente de tipo finito para toda
  extensión $S'\to S$ del preesquema de base.
  \item[\textup{(iv)}] Si $f:X\to X'$ y $g:Y\to Y'$ son dos $S$-morfismos
  localmente de tipo finito, $f\times_Sg$ es localmente de tipo finito.
  \item[\textup{(v)}] Si el compuesto $g\circ f$ de dos morfismos es localmente de
  tipo finito, $f$ es localmente de tipo finito.
  \item[\textup{(vi)}] Si un morfismo $f$ es localmente de tipo finito, lo mismo
  ocurre con $f_{\mathrm{red}}$.
\end{enumerate}
\end{proposition}

Para la demostración, véase (\textbf{I}, 6.6.6).

\begin{corollary}[1.3.5]
\label{IV.1.3.5-fr}
Sea $f:X\to Y$ un morfismo localmente de tipo finito. Para todo morfismo
$Y'\to Y$ tal que $Y'$ sea localmente noetheriano, $X\times_YY'$ es
localmente noetheriano.
\end{corollary}

\begin{proof}
En efecto, $f_{(Y')}:X\times_YY'\to Y'$ es localmente de tipo finito por
(1.3.4, (iii)), y basta aplicar (1.3.3).
\end{proof}

\begin{proposition}[1.3.6]
\label{IV.1.3.6-fr}
Sea $\varphi:A\to B$ un homomorfismo de anillos. Para que el morfismo
correspondiente $f:\operatorname{Spec}(B)\to\operatorname{Spec}(A)$ sea localmente de tipo finito, es necesario y
suficiente que $B$ sea una $A$-álgebra de tipo finito.
\end{proposition}

Para la demostración, véase (\textbf{I}, 6.3.3), teniendo en cuenta que $f$ es
cuasicompacto y (\textbf{I}, 6.6.3).

\begin{proposition}[1.3.7]
\label{IV.1.3.7-fr}
Para que un morfismo localmente de tipo finito $f:X\to Y$ sea sobreyectivo, es
necesario y suficiente que, para todo cuerpo algebraicamente cerrado $\Omega$,
la aplicación $X(\Omega)\to Y(\Omega)$ correspondiente a $f$ \textup{(\textbf{I}, 3.4.1)} sea
sobreyectiva.
\end{proposition}

Para la demostración, véase (\textbf{I}, 6.3.10).

\begin{env}[1.3.8]
\label{IV.1.3.8-fr}
Sea $A$ un anillo. Se dice que una $A$-álgebra $B$ es \emph{esencialmente de
tipo finito} si $B$ es $A$-isomorfa a una $A$-álgebra de la forma $S^{-1}C$,
donde $C$ es una $A$-álgebra \emph{de tipo finito} y $S$ un subconjunto multiplicativo
de $C$.
\end{env}

\par\medskip\noindent
\begin{proposition}[1.3.9]
\label{IV.1.3.9-fr}
\begin{enumerate}
  \item[\textup{(i)}] Si $B$ es una $A$-álgebra esencialmente de tipo finito y $C$
  una $B$-álgebra esencialmente de tipo finito, entonces $C$ es una
  $A$-álgebra esencialmente de tipo finito.
  \item[\textup{(ii)}] Sean $B$ una $A$-álgebra esencialmente de tipo finito, $A'$
  una $A$-álgebra; entonces $B'=B\otimes_AA'$ es una $A'$-álgebra
  esencialmente de tipo finito.
\end{enumerate}
\end{proposition}

\begin{proof}
\begin{enumerate}
  \item[(i)] Sean $B=S'^{-1}B'$ y $C=T'^{-1}C'$, donde $B'$ (resp. $C'$) es
  una $A$-álgebra (resp. una $B$-álgebra) de tipo finito y $S'$ (resp. $T'$)
  un subconjunto multiplicativo de $B'$ (resp. $C'$); entonces $C'$ es de la forma
  $S'^{-1}C''$, donde $C''$ es una $B'$-álgebra de tipo finito, por lo que $C$ es también
  un anillo de fracciones de $C''$; como $C''$ es una $A$-álgebra de tipo
  finito, esto prueba nuestra afirmación.
  \item[(ii)] Si $B$ es un anillo de fracciones de una $A$-álgebra de tipo finito
  $C$, $B'$ es un anillo de fracciones de la $A'$-álgebra de tipo finito
  $C\otimes_AA'$, de donde (ii).
\end{enumerate}
\end{proof}

\oldpage[IV-1]{230}

\begin{env}[1.3.10]
\label{IV.1.3.10-fr}
Si $B$ es una $A$-álgebra \emph{local} esencialmente de tipo finito, $B$ es
de la forma $C_{\mathfrak q}$, donde $C$ es una $A$-álgebra de tipo finito y
$\mathfrak q$ un ideal primo de $C$ (Bourbaki, \emph{Alg. comm.}, cap. II,
\S~5, prop. 11). Sea $\mathfrak p$ la imagen inversa en $A$ del ideal
$\mathfrak q$; si se pone $S=A\setmin\mathfrak p$, $C_{\mathfrak q}$ es también un
anillo local en un ideal primo de $S^{-1}C$; como $S^{-1}C$ es una
álgebra de tipo finito sobre $A_{\mathfrak p}=S^{-1}A$, se ve que $B$ es también
una $A_{\mathfrak p}$-\emph{álgebra esencialmente de tipo finito}, y además
el homomorfismo $A_{\mathfrak p}\to B$ es \emph{local}.
\end{env}

\begin{proposition}[1.3.11]
\label{IV.1.3.11-fr}
Si $B$ es una $A$-álgebra local esencialmente de tipo finito, $B$ es
$A$-isomorfa a una $A$-álgebra cociente de una $A$-álgebra de la forma
$C_{\mathfrak q}$, donde $C$ es un álgebra de polinomios
$A[T_1,\ldots,T_n]$ y $\mathfrak q$ un ideal primo de $C$.
\end{proposition}

\begin{proof}
En efecto, por definición, $B$ es isomorfa a $C'_{\mathfrak q'}$, donde $C'$ es
una $A$-álgebra de tipo finito y $\mathfrak q'$ un ideal primo de $C'$; pero
$C'=C/\mathfrak b$, donde $C=A[T_1,\ldots,T_n]$ y $\mathfrak b$ es un ideal de
$C$; por tanto $\mathfrak q'=\mathfrak q/\mathfrak b$, donde $\mathfrak q$ es un
ideal primo de $C$, y $C'_{\mathfrak q'}$ es isomorfa a
$C_{\mathfrak q}/\mathfrak bC_{\mathfrak q}$.
\end{proof}

